% Generated draft: requires mathematical and source review.
\documentclass[11pt]{article}
\usepackage[margin=25mm]{geometry}
\usepackage{fontspec}
\setmainfont{Latin Modern Roman}
\usepackage{amsmath,amssymb,mathtools,unicode-math}
\setmathfont{Latin Modern Math}
\usepackage{enumitem,needspace,xurl,hyperref}
\setlength{\parindent}{0pt}
\setlength{\parskip}{5pt}
\setlength{\emergencystretch}{3em}
\newcommand{\N}{\mathbb{N}}
\newcommand{\Z}{\mathbb{Z}}
\newcommand{\Q}{\mathbb{Q}}
\newcommand{\R}{\mathbb{R}}
\newcommand{\C}{\mathbb{C}}
\newcommand{\F}{\mathbb{F}}
\newcommand{\A}{\mathbb{A}}
\newcommand{\PP}{\mathbb{P}}
\newcommand{\E}{\mathbb{E}}
\newcommand{\eps}{\varepsilon}
\newcommand{\dd}{\,\mathrm{d}}
\begin{document}
\section*{UnsolvedMath ID: 701408}
\textbf{Source problem code:} catalogue entry
\subsection{Definitions and mathematical statement}
Let
\[
G=\mathrm{SL}_2(\C)=\{A\in M_2(\C):\det A=1\}.
\]
A subgroup \(H\leq G\) is free of rank two if there exist \(a,b\in H\) such that \(H=\langle a,b\rangle\) and every nonempty reduced word in \(a^{\pm1},b^{\pm1}\) represents a nonidentity element.

Does there exist a subgroup \(H\leq G\) free of rank two and an element \(g\in G\) such that
\[
gHg^{-1}\subsetneq H?
\]
Here \(gHg^{-1}=\{ghg^{-1}:h\in H\}\), and the inclusion must be proper.

The supplied source reports that a free subgroup of rank six satisfying this same conjugation property has been found by D. Calegari and N. M. Dunfield. This background assertion is retained as a source report and is not independently verified here. The registry labels the rank-two question open; that label is supplied curation rather than an independently verified assessment.
\subsection{Short English statement}
Does \(\mathrm{SL}_2(\C)\) contain a free group of rank two that conjugation by some element of \(\mathrm{SL}_2(\C)\) maps onto a proper subgroup of itself? The source reports an example of rank six.
\subsection{Sources}
\begin{itemize}
\item Local UnsolvedMath registry, permanent ID \(701408\), source problem code \texttt{catalogue entry}, title \emph{Kourovka Notebook Problem 17.107}. The supplied record labels the problem open and was updated on 2026-10-16.
\item \emph{The Kourovka Notebook (Unsolved Problems in Group Theory)}, 17th issue (2010), Problem 17.107, proposed by M. Sapir. The supplied background also identifies the 21st issue edition (2026), \texttt{arXiv:1401.0300v48}. The collection description names E. I. Khukhro and V. D. Mazurov as editors.
\item D. Calegari and N. M. Dunfield, \emph{Proc. Amer. Math. Soc.} \textbf{134}, no. 11 (2006), 3131--3136. Supplied reference for the reported rank-six example.
\end{itemize}
\end{document}
